\section{Open Questions}
Beyond the reproduced claims, five substantive follow-ups remain open:

\begin{enumerate}
\item \textbf{Does the paper's parallel-collision-finding scaling shape
hold under the actual DP-walk + collector/verifier separation, and does
the empirically optimal architecture ratio $\{1/8, 1/4, \text{rest}\}$
reproduce on modern hardware?}\\
\emph{Basis:} our \texttt{multiprocessing.Pool} implementation
reproduces the scaling \emph{shape} via first-finder-wins over
contiguous chunks, but the paper's algorithm is
van~Oorschot--Wiener DP-walks with distinct processor classes.
Fig.~4's architecture-ratio optimum is untested outside BG/Q.\\
\emph{Next steps:} implement a reference DP-walk in Python or Rust with
explicit collector/verifier processes over ZeroMQ/MPI4Py; sweep
collector-fraction $c\in\{1/16,1/8,1/4\}$ and verifier-fraction
$v\in\{1/8,1/4,1/2\}$ at $N\in\{32,64,128\}$ on uicgpu or spark-36ac;
compare best-config wall time to first-finder-wins baseline.

\item \textbf{Can pQCS-style exact $T$-count synthesis be
GPU-accelerated (FP16/BF16 tensor-core unitary comparison + massive
batched DP-walks on a single 8-GPU box) to close the gap with a
4096-core BG/Q run?}\\
\emph{Basis:} the per-worker work is unitary multiplication + hashing,
both trivially data-parallel and both GPU-friendly at $8{\times}8$ or
$16{\times}16$ scale. No published pQCS variant runs on GPUs.\\
\emph{Next steps:} port the DP-walk inner loop to CuPy or JAX \texttt{pmap}
on the 8 A100s at uicgpu; batch-evaluate $10^6$ walks per launch;
profile hash vs matmul overhead; benchmark against the paper's 26\,s
Toffoli figure.

\item \textbf{Does a hybrid classical/quantum pipeline (variational
Clifford+$T$ ansatz on IBM/IonQ hardware for coarse decomposition, then
pQCS for exact refinement of sub-blocks) beat either endpoint alone on
5--8 qubit targets?}\\
\emph{Basis:} pQCS scales exponentially in depth; VQC/VQE hardware is
noisy but samples a continuous ansatz space cheaply. No hybrid split has
been demonstrated for exact $T$-count-optimal synthesis.\\
\emph{Next steps:} take a 6-qubit target (small qiskit-nature chemistry
unitary or QFT variant); run variational Clifford+$T$ ansatz on IBM
hardware via \texttt{qiskit-ibm-runtime}; refine each Clifford+$T$
sub-block with pQCS-style search; compare vs pQCS-only.

\item \textbf{How does pQCS compare head-to-head against modern
optimizing compilers (tket \texttt{FullPeepholeOptimise},
qiskit-transpiler \texttt{optimization\_level=3}, both of which now
include $T$-count-aware passes derived from post-2016 work)?}\\
\emph{Basis:} pQCS is 2016; tket and modern qiskit-transpiler
incorporate many post-pQCS $T$-count reduction techniques. No
head-to-head benchmark on a common target set is published.\\
\emph{Next steps:} 20 target circuits (10 Toffoli variants + 10
arithmetic benchmarks); run tket + qiskit-transpiler + pQCS reference;
tabulate final $T$-count and wall time to optimum.

\item \textbf{Can pQCS be extended to noise-aware synthesis --- picking
a $T$-count-optimal decomposition among the many equal-$T$-count
optima that minimizes expected error under a specific device noise
model --- and does that materially improve end-to-end circuit
fidelity?}\\
\emph{Basis:} pQCS finds an optimum but is agnostic among equal-cost
optima. On real hardware, equal-$T$-count decompositions differ in
depth, layout, and crosstalk exposure. No noise-aware pQCS variant is
published.\\
\emph{Next steps:} enumerate all known $T$-count-7 Toffoli
decompositions; simulate each under a realistic IBM Heron or IonQ Forte
noise model via qiskit-aer + device backend snapshot; add a weighted
expected-error penalty to the pQCS verifier; validate on IBM Quantum's
open plan.
\end{enumerate}
